\part{Nullstellen und Ideale}

\chapter{Was eine Gleichung zeichnet}

\section{Die Lösungsmenge als Gegenstand}

\begin{merksatz}
\boxtitel{Drei Gleichungen, vier Unbekannte, eine Kurve}
Für den Leitfall mit gemessenen Werten $s_1 = \tfrac12$, $s_2 = \tfrac{9}{10}$, $s_3 = \tfrac{14}{100}$ lauten die Gleichungen
\[ k_{10} + k_{12} = \tfrac12, \qquad k_{10} + k_{12} + k_{21} + k_{20} = \tfrac{9}{10}, \qquad k_{10}k_{21} + k_{10}k_{20} + k_{12}k_{20} = \tfrac{14}{100} . \]
Aus den ersten beiden: $k_{10} = \tfrac12 - k_{12}$ und $k_{21} = \tfrac{2}{5} - k_{20}$. Einsetzen in die dritte:
\[ \bigl(\tfrac12 - k_{12}\bigr)\bigl(\tfrac25 - k_{20}\bigr) + \bigl(\tfrac12 - k_{12}\bigr)k_{20} + k_{12}k_{20} = \tfrac15 - \tfrac25 k_{12} + k_{12}k_{20} = \tfrac{14}{100}, \]
also $k_{12}k_{20} - \tfrac25 k_{12} + \tfrac{3}{50} = 0$, oder $k_{20} = \tfrac25 - \dfrac{3}{50\,k_{12}}$.

\begin{center}
{\small
\renewcommand{\arraystretch}{1.46}
\begin{tabular}{@{}P{1.4cm}P{1.4cm}P{1.4cm}P{1.4cm}P{2.6cm}@{}}
\toprule
$k_{12}$ & $k_{10}$ & $k_{20}$ & $k_{21}$ & \textbf{Deutung}\\
\midrule
$0{,}15$ & $0{,}35$ & $0{,}00$ & $0{,}40$ & keine Ausscheidung aus dem Gewebe\\
$0{,}20$ & $0{,}30$ & $0{,}10$ & $0{,}30$ & die \enquote{wahren} Werte\\
$0{,}30$ & $0{,}20$ & $0{,}20$ & $0{,}20$ & gleiche Ausscheidung beiderseits\\
$0{,}50$ & $0{,}00$ & $0{,}28$ & $0{,}12$ & keine Ausscheidung aus dem Blut\\
\bottomrule
\end{tabular}}
\end{center}

\begin{center}
\begin{minipage}{0.9\textwidth}
\textbf{Jede Zeile ist ein Modell, und alle vier erzeugen exakt dieselbe Blutkurve. Die Lösungsmenge ist eine Kurve im vierdimensionalen Parameterraum --- ein Hyperbelast, parametrisiert durch $k_{12}\in[0{,}15;\ 0{,}5]$. Die Daten sagen nicht, wo auf ihr das Medikament liegt.}
\end{minipage}
\end{center}

\begin{center}
\begin{minipage}{0.88\textwidth}
\kursiv{\textbf{Das Rechnen war hier leicht, weil zwei der drei Gleichungen linear waren.} Mit drei Kammern, oder mit einer Messung im Gewebe statt im Blut, oder mit nichtlinearer Ausscheidung sind alle Gleichungen vom Grad zwei oder höher, und das Einsetzen führt in ein Dickicht. Dieses Buch ersetzt das Einsetzen durch eine Theorie: Die Lösungsmenge wird zum Gegenstand, mit Dimension, Bestandteilen und Gleichungen, die man systematisch vereinfachen kann. Die Kurve des Leitfalls hat Dimension eins --- das ist der Satz \enquote{nicht identifizierbar} in der Sprache dieses Bandes, und Kapitel 7 sagt, wie man die Dimension berechnet, ohne die Kurve zu finden.}
\end{minipage}
\end{center}
\end{merksatz}

\section{Affine Varietäten}

\begin{definition}
\boxtitelm{Definition 1.1 --- affine Varietät}
Sei $k$ ein Körper und $k[x_1, \ldots, x_n]$ der Ring der Polynome in $n$ Veränderlichen. Für Polynome $f_1, \ldots, f_s$ heißt
\[ V(f_1, \ldots, f_s) = \{\,a \in k^{n} : f_1(a) = \cdots = f_s(a) = 0\,\} \]
die von ihnen bestimmte \Begriff{affine Varietät}. Der Raum $k^{n}$ heißt \Begriff{affiner Raum} $\mathbb A^{n}$; der Körper ist in diesem Band meist $\mathbb Q$, $\mathbb R$ oder $\mathbb C$, in Kapitel 10 und 14 auch ein endlicher Körper $\mathbb F_p$.
\end{definition}

\begin{merksatz}
\boxtitel{Beispiele, die wiederkommen}
\begin{center}
{\small
\renewcommand{\arraystretch}{1.5}
\begin{tabular}{@{}P{3.6cm}P{5.0cm}@{}}
\toprule
$V(x^{2} + y^{2} - 1)\subset\mathbb R^{2}$ & der Einheitskreis: eine Kurve\\
$V(y - x^{2})$ & die Parabel\\
$V(y^{2} - x^{3} + 2x)$ & eine Kubik mit zwei Ästen (Titelbild)\\
$V(y^{2} - x^{3})$ & die Neilsche Parabel mit Spitze im Ursprung\\
$V(y^{2} - x^{2}(x+1))$ & ein Blatt, das sich im Ursprung kreuzt\\
$V(y - x^{2},\ z - x^{3})\subset\mathbb R^{3}$ & die verdrehte Kubik $(t, t^{2}, t^{3})$: eine Raumkurve\\
$V(xz,\ yz)$ & die $xy$-Ebene vereinigt mit der $z$-Achse\\
$V(x^{2} + y^{2})\subset\mathbb R^{2}$ & ein Punkt --- aber in $\mathbb C^{2}$ zwei Geraden\\
$V(1)$, $V(0)$ & leer; der ganze Raum\\
\bottomrule
\end{tabular}}
\end{center}

\begin{center}
\begin{minipage}{0.9\textwidth}
\textbf{Eine Gleichung in der Ebene gibt gewöhnlich eine Kurve, zwei Gleichungen im Raum eine Kurve, eine Gleichung im Raum eine Fläche. Aber nicht immer: $V(xz, yz)$ hat Teile verschiedener Dimension, und $V(x^{2}+y^{2})$ hängt vom Körper ab.}
\end{minipage}
\end{center}
\end{merksatz}

\begin{merksatz}
\boxtitel{Satz 1.2 --- Varietäten lassen sich schneiden und vereinigen}
Sind $V = V(f_1, \ldots, f_s)$ und $W = V(g_1, \ldots, g_t)$ affine Varietäten in $k^{n}$, so auch
\[ V\cap W = V(f_1, \ldots, f_s, g_1, \ldots, g_t) \qquad\text{und}\qquad V\cup W = V(f_ig_j : 1\le i\le s,\ 1\le j\le t). \]

\schritt{Beweis für den Durchschnitt}
$a\in V\cap W$ genau dann, wenn alle $f_i(a) = 0$ und alle $g_j(a) = 0$ --- das ist die Definition der rechten Seite. $\blacksquare$

\schritt{Beweis für die Vereinigung}
Ist $a\in V$, so sind alle $f_i(a) = 0$, also alle Produkte $f_i(a)g_j(a) = 0$: $a$ liegt rechts. Ebenso für $a\in W$. Umgekehrt: Liegt $a$ rechts und nicht in $V$, so ist ein $f_{i_0}(a)\ne0$; aus $f_{i_0}(a)g_j(a) = 0$ für alle $j$ folgt dann $g_j(a) = 0$ für alle $j$, also $a\in W$. $\blacksquare$

\schritt{Beispiel}
$V(xz, yz) = V(x, y)\cup V(z)$: die $z$-Achse ($x = y = 0$) vereinigt mit der $xy$-Ebene ($z = 0$). Die Produkte $x\cdot z$ und $y\cdot z$ sind genau die Gleichungen. \checkmark

\begin{center}
\begin{minipage}{0.88\textwidth}
\kursiv{\textbf{Der Satz sagt, dass die Varietäten die abgeschlossenen Mengen einer Topologie bilden} --- der Zariski-Topologie, Band 8 ---, in der endliche Vereinigungen und beliebige Durchschnitte abgeschlossener Mengen abgeschlossen sind. Es ist eine seltsame Topologie: Jede nichtleere offene Menge der Geraden ist das Komplement endlich vieler Punkte, zwei offene Mengen schneiden sich immer. Aber sie ist die richtige für Polynome, weil sie nur das sieht, was Polynome sehen. Die Vereinigung in Satz 1.2 ist der Grund, warum Varietäten in Bestandteile zerfallen (Kapitel 3): $V(xz, yz)$ ist nicht \enquote{ein} Gebilde, sondern zwei.}
\end{minipage}
\end{center}
\end{merksatz}

\section{Parametrisierungen}

\begin{merksatz}
\boxtitel{Satz 1.3 --- von der Parametrisierung zur Gleichung, und zurück}
Zwei Arten, eine Varietät zu beschreiben:
\begin{center}
{\small
\renewcommand{\arraystretch}{1.5}
\begin{tabular}{@{}P{2.4cm}P{3.0cm}P{3.2cm}@{}}
\toprule
 & \textbf{implizit} & \textbf{parametrisch}\\
\midrule
Form & $f_1 = \cdots = f_s = 0$ & $x_i = g_i(t_1, \ldots, t_m)$\\
gut für & prüfen, ob ein Punkt dazugehört & Punkte erzeugen, zeichnen\\
Beispiel Kreis & $x^{2} + y^{2} = 1$ & $x = \frac{1-t^{2}}{1+t^{2}},\ y = \frac{2t}{1+t^{2}}$\\
\bottomrule
\end{tabular}}
\end{center}
Jede polynomiale (oder rationale) Parametrisierung beschreibt eine Teilmenge einer affinen Varietät, und die kleinste Varietät, die sie enthält, lässt sich berechnen (\Begriff{Implizitisierung}, Kapitel 6). Nicht jede Varietät hat eine Parametrisierung (Kapitel 12).

\schritt{Der Kreis, nachgerechnet}
$x^{2} + y^{2} = \dfrac{(1-t^{2})^{2} + 4t^{2}}{(1+t^{2})^{2}} = \dfrac{1 + 2t^{2} + t^{4}}{(1+t^{2})^{2}} = 1$. \checkmark\ Die Parametrisierung trifft jeden Punkt des Kreises außer $(-1, 0)$: Der wäre $t = \infty$.

\schritt{Die verdrehte Kubik}
$(t, t^{2}, t^{3})$ erfüllt $y = x^{2}$ und $z = x^{3}$; umgekehrt gibt jeder Punkt von $V(y - x^{2}, z - x^{3})$ mit $t = x$ eine Parameterdarstellung. Hier stimmen Bild und Varietät überein.

\begin{center}
\begin{minipage}{0.9\textwidth}
\textbf{Der Leitfall ist implizit gegeben (drei Gleichungen) und wurde oben parametrisiert ($k_{12}$ als Parameter). Die Richtung von implizit nach parametrisch ist im Allgemeinen schwer oder unmöglich; die Richtung von parametrisch nach implizit ist immer möglich und ist eine Elimination.}
\end{minipage}
\end{center}

\begin{center}
\begin{minipage}{0.88\textwidth}
\kursiv{\textbf{Die Parametrisierung des Kreises ist die von Diophant und Euklid}: Setzt man $t = p/q$, erhält man die pythagoreischen Tripel $(q^{2} - p^{2},\ 2pq,\ q^{2} + p^{2})$ --- jedes rationale Punktepaar auf dem Kreis aus einer rationalen Zahl $t$, Band 1. Dass der Kreis eine solche Parametrisierung hat und die Kubik des Titelbilds nicht, ist der Unterschied zwischen Geschlecht null und Geschlecht eins (Kapitel 12), und er entscheidet darüber, ob eine Gleichung unendlich viele rationale Lösungen hat oder endlich viele: Fermats Vermutung ist eine Aussage über das Geschlecht der Kurve $x^{n} + y^{n} = 1$.}
\end{minipage}
\end{center}
\end{merksatz}

\section{Der Bruchfall}

\begin{warnung}
\boxtitelr{Die Zeichnung zeigt die Varietät nicht}
Man zeichnet $V(x^{2} + y^{2})$ in der reellen Ebene: ein Punkt. Man zeichnet $V(x^{2} + y^{2} + 1)$: nichts. Man zeichnet $V(y^{2} - x^{2}(x - 1))$: eine Kurve rechts von $x = 1$ --- und einen einsamen Punkt im Ursprung.

\begin{center}
{\small
\renewcommand{\arraystretch}{1.5}
\begin{tabular}{@{}P{3.0cm}P{2.6cm}P{2.8cm}@{}}
\toprule
\textbf{Gleichung} & \textbf{über $\mathbb R$} & \textbf{über $\mathbb C$}\\
\midrule
$x^{2} + y^{2} = 0$ & ein Punkt & zwei Geraden $y = \pm\mathrm ix$\\
$x^{2} + y^{2} = -1$ & leer & eine Kurve (ein Kegelschnitt)\\
$y^{2} = x^{2}(x-1)$ & Kurve und isolierter Punkt & eine zusammenhängende Kurve\\
$x^{2} = 0$ & die $y$-Achse & die $y$-Achse --- doppelt\\
\bottomrule
\end{tabular}}
\end{center}

\begin{center}
\begin{minipage}{0.9\textwidth}
\textbf{Über den reellen Zahlen kann ein Polynom vom Grad zwei eine leere Nullstellenmenge haben, eine Varietät kann Dimension null und zwei zugleich zeigen, und zwei verschiedene Gleichungen können dasselbe Bild haben. Keines davon geschieht über $\mathbb C$.}
\end{minipage}
\end{center}

\begin{center}
\begin{minipage}{0.88\textwidth}
\kursiv{\textbf{Deshalb arbeitet die algebraische Geometrie über einem algebraisch abgeschlossenen Körper} --- über $\mathbb C$, in dem jedes Polynom seine Nullstellen hat (Band 12). Dort gilt: Jede nichtkonstante Gleichung hat Lösungen, die Dimension ist überall gleich auf einem Bestandteil, und die Gleichungen bestimmen die Varietät bis auf Vielfachheit. Der isolierte Punkt der dritten Zeile ist in $\mathbb C^{2}$ ein Punkt, in dem die Kurve sich selbst kreuzt --- mit zwei komplexen Ästen, die man reell nicht sieht. Für den Leitfall ist der Körper $\mathbb R$, und mehr noch: die Parameter müssen positiv sein. Die Theorie wird über $\mathbb C$ gemacht, und am Ende fragt man, welche der komplexen Lösungen reell und positiv sind; Kapitel 13 zeigt, dass diese letzte Frage eine eigene ist.}
\end{minipage}
\end{center}
\end{warnung}

\section{Übungsaufgaben}

\begin{aufgabe}
\boxtitel{Aufgabe 1.1 --- Zerlegen}
Beschreiben Sie $V(x^{2} - y^{2},\ xz - yz)\subset\mathbb R^{3}$ als Vereinigung einfacher Varietäten.

\loesung
\schritt{Schritt 1 --- Faktorisieren}
$x^{2} - y^{2} = (x-y)(x+y)$ und $xz - yz = (x-y)z$.

\schritt{Schritt 2 --- Fallunterscheidung}
Ist $x = y$, so sind beide Gleichungen erfüllt, für jedes $z$: die Ebene $V(x - y)$. Ist $x\ne y$, so muss $x + y = 0$ und $z = 0$ sein: die Gerade $V(x + y, z)$.

\schritt{Schritt 3 --- Ergebnis}
$V = V(x - y)\cup V(x + y,\ z)$: eine Ebene und eine Gerade, die in ihr nicht liegt (die Gerade $x = -y$, $z = 0$ trifft die Ebene $x = y$ nur im Ursprung).

\schritt{Schritt 4 --- Probe mit Satz 1.2}
$V(x-y)\cup V(x+y, z) = V\bigl((x-y)(x+y),\ (x-y)z\bigr) = V(x^{2}-y^{2},\ xz-yz)$. \checkmark

\begin{center}
\begin{minipage}{0.88\textwidth}
\kursiv{\textbf{Die Varietät hat Teile der Dimension zwei und eins.} Keiner lässt sich weiter zerlegen --- eine Ebene und eine Gerade sind \enquote{irreduzibel}, Kapitel 3 ---, und die Zerlegung ist eindeutig. Dass man sie hier durch Faktorisieren fand, ist Glück; im Allgemeinen braucht es die Primzerlegung des Ideals, Kapitel 3 und 5.}
\end{minipage}
\end{center}
\end{aufgabe}

\begin{aufgabe}
\boxtitel{Aufgabe 1.2 --- Der Leitfall mit drei Parametern}
Das Modell ohne Ausscheidung aus dem Gewebe ($k_{20} = 0$) liefert $s_1 = k_{10} + k_{12} = \tfrac12$, $s_2 = k_{10} + k_{12} + k_{21} = \tfrac45$, $s_3 = k_{10}k_{21} = \tfrac{9}{100}$. Bestimmen Sie die Varietät.

\loesung
\schritt{Schritt 1}
Aus $s_2 - s_1$: $k_{21} = \tfrac{3}{10}$.

\schritt{Schritt 2}
Aus $s_3$: $k_{10} = \tfrac{9}{100}\big/\tfrac{3}{10} = \tfrac{3}{10}$.

\schritt{Schritt 3}
Aus $s_1$: $k_{12} = \tfrac12 - \tfrac{3}{10} = \tfrac15$.

\schritt{Schritt 4 --- Probe}
$k_{10} + k_{12} = 0{,}5$, $+ k_{21} = 0{,}8$, $k_{10}k_{21} = 0{,}09$. \checkmark\ Die Varietät ist ein einziger Punkt $(0{,}3;\ 0{,}2;\ 0{,}3)$: Das Modell ist \textbf{global identifizierbar}.

\begin{center}
\begin{minipage}{0.88\textwidth}
\kursiv{\textbf{Ein Parameter weniger, und aus der Kurve wird ein Punkt.} Das ist die Grundregel der Identifizierbarkeit: Drei Datenzahlen bestimmen höchstens drei Parameter --- aber nicht immer drei (Aufgabe 1.3), und welche drei, sagt erst die Rechnung. Der Pharmakologe weiß meist aus anderen Gründen, dass $k_{20} = 0$ (die Niere hängt am Blut, nicht am Gewebe); die Algebra sagt ihm, dass er dieses Wissen braucht.}
\end{minipage}
\end{center}
\end{aufgabe}

\begin{aufgabe}
\boxtitel{Aufgabe 1.3 --- Messung in der falschen Kammer}
Misst man statt im Blut im Gewebe (Kammer 2) und setzt $k_{20} = 0$, so liefert die Messung $k_{12} = \tfrac15$, $k_{10} + k_{12} + k_{21} = \tfrac45$ und $k_{10}k_{21} = \tfrac{8}{100}$. Bestimmen Sie die Varietät.

\loesung
\schritt{Schritt 1}
$k_{12} = \tfrac15$ ist direkt gegeben; also $k_{10} + k_{21} = \tfrac35$ und $k_{10}k_{21} = \tfrac{8}{100}$.

\schritt{Schritt 2 --- Vieta}
$k_{10}$ und $k_{21}$ sind die Nullstellen von $u^{2} - \tfrac35u + \tfrac{8}{100} = 0$: $u = \tfrac{3}{10}\pm\sqrt{\tfrac{9}{100} - \tfrac{8}{100}} = \tfrac{3}{10}\pm\tfrac{1}{10}$, also $u = \tfrac25$ oder $\tfrac15$.

\schritt{Schritt 3 --- Zwei Punkte}
$(k_{10}, k_{12}, k_{21}) = (0{,}4;\ 0{,}2;\ 0{,}2)$ oder $(0{,}2;\ 0{,}2;\ 0{,}4)$. Probe: $0{,}4 + 0{,}2 + 0{,}2 = 0{,}8$, $0{,}4\cdot0{,}2 = 0{,}08$. \checkmark\ Beide Male.

\schritt{Schritt 4 --- Deutung}
Die Varietät besteht aus zwei Punkten: \textbf{lokal identifizierbar} (um jede Lösung herum ist sie isoliert), aber nicht global --- die Daten unterscheiden nicht, ob der Wirkstoff schnell ausgeschieden und langsam zurückverteilt wird oder umgekehrt.

\begin{center}
\begin{minipage}{0.88\textwidth}
\kursiv{\textbf{Punkt, Punktepaar, Kurve --- drei Modelle, drei Gestalten.} Die Zahl der Punkte einer nulldimensionalen Varietät ist ihr Grad (Kapitel 7 und 9); hier ist er zwei, weil die Gleichung in $u$ quadratisch ist. Für größere Modelle ist der Grad der Fasern eine der wichtigsten Zahlen der Identifizierbarkeitstheorie: Er sagt, wie viele Modelle man gegeneinander abwägen muss, nachdem man gemessen hat.}
\end{minipage}
\end{center}
\end{aufgabe}

\begin{geschichte}
\boxtitelg{Descartes' Koordinaten und Newtons Kurven}
\Person{René Descartes} veröffentlichte 1637 als Anhang zum \enquote{Discours de la méthode} die \enquote{Géométrie}: Kurven werden durch Gleichungen beschrieben, und geometrische Probleme werden durch Rechnen gelöst. Die Gleichung war fortan die Kurve.

\begin{center}
\begin{minipage}{0.9\textwidth}
\kursiv{\Person{Fermat} hatte dasselbe gleichzeitig und unabhängig gefunden, veröffentlichte aber nicht. \textbf{\Person{Newton} klassifizierte um 1670 die ebenen Kurven dritten Grades}: $72$ Arten, später auf $78$ ergänzt, nach Gestalt und Lage im Unendlichen --- die erste systematische Untersuchung einer Varietätenklasse. Er bemerkte, dass alle aus fünf Grundformen durch Projektion hervorgehen; das war die projektive Geometrie, Kapitel 9, hundertfünfzig Jahre vor ihrem Namen.}
\end{minipage}
\end{center}

\schritt{Das Wort}
\begin{center}
\begin{minipage}{0.9\textwidth}
\kursiv{\enquote{Varietät} kommt von \Person{Riemann}s \enquote{Mannigfaltigkeit} (1854) über die italienische Schule --- \Person{Castelnuovo}, \Person{Enriques}, \Person{Severi} ---, die um 1900 Flächen klassifizierte und \enquote{varietà} sagte. \textbf{Die Definition 1.1 in dieser Form ist jung}: Sie setzt voraus, dass man Polynomringe als Gegenstand kennt (\Person{Hilbert} 1890, Kapitel 2) und dass man weiß, welcher Körper gemeint ist --- eine Frage, die erst \Person{Weil} 1946 für beliebige Körper beantwortete.}
\end{minipage}
\end{center}
\end{geschichte}

\begin{forschung}
\boxtitelp{Identifizierbarkeit als Geometrie}
\begin{center}
{\small
\renewcommand{\arraystretch}{1.45}
\begin{tabular}{@{}P{3.0cm}P{5.2cm}@{}}
\toprule
\textbf{Strukturelle Identifizierbarkeit} & Für Kompartimentmodelle mit Dutzenden Parametern: Programme (DAISY, SIAN, StructuralIdentifiability.jl) berechnen, welche Parameter aus welchen Messungen folgen --- mit den Methoden der Kapitel 5 und 6.\\
\textbf{Praktische Identifizierbarkeit} & Ein Punkt in der Theorie, eine flache Mulde in den Daten: Wenn die Varietät ein Punkt ist, aber die Daten verrauscht, ist der Parameter zwar bestimmt, aber schlecht --- Band 17 trifft Band 29.\\
\textbf{Modellreduktion} & Wenn die Varietät eine Kurve ist: Welche Funktion der Parameter ist konstant auf ihr? Das ist der identifizierbare Anteil, und er ist ein Element des Koordinatenrings (Kapitel 3).\\
\bottomrule
\end{tabular}}
\end{center}

\schritt{Zur dritten Zeile}
\begin{center}
\begin{minipage}{0.88\textwidth}
\kursiv{\textbf{Auf der Kurve des Leitfalls sind $k_{10} + k_{12}$, $k_{21} + k_{20}$ und $k_{12}k_{20} - \tfrac25k_{12}$ konstant} --- das sind die identifizierbaren Kombinationen, und ein Arzt, der ein Medikament dosiert, braucht oft nur sie: Die Gesamtausscheidung aus dem Blut, $k_{10} + k_{12}$, bestimmt die Anfangsphase, und die Summe $k_{21} + k_{20}$ die Rückkehr. Die Frage, welche Funktionen auf einer Varietät konstant sind, ist die Frage nach dem Koordinatenring, Kapitel 3, und die Antwort ist in diesem Fall berechenbar: Es sind genau die Polynome in $s_1, s_2, s_3$. Band 26 hat dieselbe Frage für lineare Modelle gestellt; hier ist sie für polynomiale beantwortbar.}
\end{minipage}
\end{center}
\end{forschung}

\begin{merksatz}
\boxtitel{Was Sie jetzt können}
\begin{center}
\renewcommand{\arraystretch}{1.5}
\begin{tabular}{@{}P{3.0cm}P{5.6cm}@{}}
\toprule
\textbf{Definition 1.1} & affine Varietät $V(f_1, \ldots, f_s)$\\
\textbf{Satz 1.2} & Durchschnitt und Vereinigung sind Varietäten\\
\textbf{Satz 1.3} & implizit und parametrisch\\
\textbf{Der Bruchfall} & reell gezeichnet ist nicht die Varietät; arbeite über $\mathbb C$\\
\textbf{Leitfall} & Punkt, Punktepaar, Kurve = global, lokal, nicht identifizierbar\\
\bottomrule
\end{tabular}
\end{center}

\begin{center}
\begin{minipage}{0.88\textwidth}
\kursiv{Die Varietät ist der Gegenstand. Kapitel 2 fragt nach den Gleichungen: Welche Polynome verschwinden auf einer Varietät? Die Antwort ist ein Ideal --- und Hilberts Satz, dass es endlich erzeugt ist.}
\end{minipage}
\end{center}
\end{merksatz}
